\documentclass[a4paper, oneside]{book}
\usepackage[italian]{babel}
\usepackage{graphicx} % allows inserting graphics
\usepackage{listings}
\usepackage{algpseudocode} % allows code snippets
\usepackage{algorithm}
\usepackage{hyperref} % allows
\usepackage{subcaption} 
\usepackage[export]{adjustbox} 
\usepackage{wrapfig} % allows warped figures
\usepackage[acronym]{glossaries} % add glossary
\usepackage{acronym} % add acronyms
\usepackage{subfiles} % support multiple files
\usepackage{fancyhdr} % more control on headers and footers
\usepackage{xcolor}
\usepackage{threeparttable}

\usepackage{arydshln}
\usepackage{multirow}
\usepackage{titlesec}

\newcommand{\red}[1]{\textcolor{red}{#1}}

\graphicspath{ {./images/} }

\pagestyle{plain} %oppure none
% Altrimenti usare fancy header per più controllo: https://www.overleaf.com/learn/latex/Headers_and_footers#Using_the_fancyhdr_package

\usepackage[
backend=biber,
style=ieee,
sorting=ynt
]{biblatex}

% serve per risolvere conflitti di versione della libreria biblatex
\makeatletter
\@ifundefined{abx@macro@volume+number}{%
  \newbibmacro*{volume+number}{%
    \printfield{volume}%
    \setunit*{\adddot}%
    \printfield{number}%
  }%
}{}
\makeatother

\addbibresource{extra/bibliografia.bib}
% per mostrare tutta la bibliografia, anche ciò che non è stato citato
% \nocite{*} 

\makeglossaries
\loadglsentries{extra/glossario.tex}
\setcounter{tocdepth}{2}

\begin{document}

\subfile{chapters/0-frontespizio.tex}

\tableofcontents
% \listoffigures
% \listoftables
\printglossary
\clearpage

% \subfile{extra/acronimi.tex}

\subfile{chapters/1-introduzione.tex}
\subfile{chapters/2-stato-dell-arte.tex}
\subfile{chapters/3-dataset.tex}
% \subfile{chapters/4-metodi.tex}
\subfile{chapters/5-esperimenti-e-risultati.tex}
\subfile{chapters/6-conclusioni.tex}

% \subfile{chapters/tutorial/1-structure.tex}
% \subfile{chapters/tutorial/2-text-style.tex}
% \subfile{chapters/tutorial/3-elements-of-the-document.tex}
% \subfile{chapters/tutorial/4-positioning-and-size.tex}
% \subfile{chapters/tutorial/5-advanced-images-and-table.tex}
% \subfile{chapters/tutorial/6-additional-concepts.tex}
% \subfile{chapters/tutorial/7-multiple-files-management.tex}
% \subfile{chapters/tutorial/8-citations.tex}

\subfile{chapters/7-bibliografia.tex}

% \appendix
% \subfile{chapters/8-appendice-a.tex}
% \subfile{chapters/9-appendice-b.tex}

\end{document}